\paragraph{\trGrading{}问题。}\label{sub:grading-issues}
设 $M^{2n}$ 是一个\trCorneredBoundary{}的\trExactSymplectic\trManifold，其上带有二次复\trVolume{}形式 $\eta^2_M$。若 $V\subset M$ 是一个容许\trGrading{}的带\trFrame{}的\trExact\trLagrange\trSphere（无须选定任何特定的\trGrading），那么\trDehnTwist{}$\tau_V$ 就自然是 \cref{sub:automorphisms} 意义下的一个\trGraded\trSymplectic\trAutomorphism；实值\trGrading{}函数 $\gamma^{\#}_{\tau_V}$ 由它在 $V$ 的某个邻域之外恒为零这一条件确定（在这个邻域中 $\tau_V$ 是恒等的）。其证明并不困难：一次\trCohomology{}检验表明，在维数 $n>1$ 时没有\trObstruction，而剩下的情形是初等的。
另一方面，考虑一个带有相对二次复\trVolume{}形式 $\eta^2_{E/S}$ 的\trLefschetzFibration。它在诸\trSmooth\trFiber{}上诱导出一些形式；相对于这些形式，任何\trVanishingCycle{}都容许一个\trGrading{}（但并不存在哪一种选择是自然特异的）。本来也可以像前面那样论证——即除最低维的那个情形可由手工解决之外，其余情形都是平凡的。然而，还有一个更为漂亮的论证，它与证明\trVanishingCycle{}\trExact{}性的论证相类似：设 $\gamma$ 为\trVanishingPath。在 $\pi^{-1}(\gamma)\subset E$ 的一个邻域内，写 $\eta^2_{E/S}=\eta^2_E/\eta^2_S$（其实处处都可以这样写，因为由假定 $TS$ 是平凡的，不过我们并不需要这一点）。由于 $\Delta_\gamma$ 是\trContractible{}的，与 $\eta^2_E$ 相伴的\trSquaredPhase{}函数
\begin{equation*}
\alpha_{\Delta_\gamma}\colon\Delta_\gamma\longrightarrow S^1
\end{equation*}
当然容许一个到 $\RR$ 的\trLift。该函数限制到边界上时，与 $V_\gamma$ 的\trSquaredPhase{}函数相差一个常数（这个常数依赖于 $\gamma'(0)$），于是所要的事实便随之而得。最后，查看上面所概述的 \eqref{eq:symplectic-picard-lefschetz} 的证明即可看出：如果给\trMonodromy{}配上由 \eqref{eq:relative-grading-function} 所定义的\trGrading，并按我们刚才描述的方式为\trDehnTwist{}\trGrading，那么那里构造的\trIsotopy{}便可\trLift{}到由\trGraded\trSymplectic\trAutomorphism{}组成的群 $\mathrm{Aut}^{\mathrm{gr}}(M)$ 之中。
最后，如果给定 $M$、一个 $\eta^2_M$，以及一族都容许\trGrading{}的带\trFrame{}的\trExact\trLagrange\trSphere，那么按\cref{thm:realizing-distinguished-basis}构造出来的那个相伴\trLefschetzFibration{}将容许一个相对二次复\trVolume{}形式，它在相对于 $E_*=M$ 的\trHomotopy{}意义下唯一。同样，这里也没有理由去作一个特定的\trGrading{}选择（要紧的只是存在性）。
\begin{remark}\label{thm:dehn-twist-history-remarks}
高维\trDehnTwist{}是辛的这一事实，系\trArnold{}所发现 [3]。至于\trVanishingCycle{}相应的\trLagrange{}性质，则是\trDonaldson{}向作者指出的，尽管这很可能早为人知。至于\trPicard--\trLefschetz{}定理，它在\trTopology{}的层面上（把\trMonodromy{}当作\trDiffeomorphism）乃是经典的。众多可能的文献中，可举 [92] 与 [5] 两种。辛的版本可在 [131, Proposition 1.15] 中找到。\trVanishingCycle{}的\trDistinguishedBasis{}同样是一个经典概念，在\trSingularPoint{}理论 [5]、[4] 所考虑的那些局部情形中尤其如此。“在辛的层面上，\trTotalSpace{}可以由这样一个基重构出来”这一观察则较为晚近（参见本章开头所引的诸篇论文）。\trDehnTwist{}的\trGrading{}在 [125] 中有更详细的讨论。
\end{remark}
